\scanpage{115}
% archive foliation: 116
\overline{\underline{\mathrm{III}}}) Dans le cas où $S=\operatorname{Spec}(k)$, $k$ corps parfait de car. $p>0$, on retrouve le foncteur de Dieudonné habituel.

\overline{\underline{\mathrm{IV}}}) (Déformations) $D(G)_S$ est filtré par
\[
0\longrightarrow\overset{\deleted{\unsure{\vee}}}{\unsure{\omega_{G^*}}}\longrightarrow D(G)_S\longrightarrow\omega_{G^*}^{\vee}\longrightarrow0.
\]
Soit $S_o\subset S$
\margin{ép. à p. div.}
Supposons \deleted{\unsure{que}} \inserted{\unsure{que les p. div. sur $I$ sont nilpotentes}}, \unsure{ou} bien \unsure{aussi exige} $G_o$ tel que les fibres n'aient pas de comp. toroïdale, ou que les fibres de $G_o$ soient connexes. Partant de $G_o$, on trouve que prolonger $G_o$, c'est ``la $=$ chose'' que de prolonger la filtration de $D(G_o)_{S_o}$ à $D(G_o)_S$.

\overline{\underline{\mathrm{V}}}) Prenons les groupes de BT à isog. près.

Supposons $A$ séparé et complet pour la top. $(pA)$-adique ; $A_o=A/pA$. Alors foncteur
\[
G\mathrel{\unsure{\mapstochar\relbar\relbar|\longmapsto}}(G_o,\text{filtration sur }D(G_o)_A)
\]
de
\[
\mathrm{BT}_{\mathrm{isog}}(S)\longrightarrow\unsure{\mathcal C--}
\]
est pleinement fidèle.

\uline{Cas d'un anneau} \deleted{de valuation \unsure{discrète}} \inserted{local complet \unsure{noeth}, à corps résiduel} \uline{$k$ parfait de car.} $p>0$ : Groupe de B.T. est connu \inserted{(mod. isogénie)} quand on connaît $G_o$ \uline{sur} $k$ (i.e. module de Dieud. $M$ sur $k$) et une filtration sur $(M\otimes_W A)_{\unsure{p|\mathbb Q}}$
